\chapter{Números enteros}\label{ch:g6:wholes}

Todo en matemáticas comienza con contar. Este capítulo analiza cómo
números enteros se escriben, comparan y combinan --- y toma una cuidadosa
Observemos la división, la operación que da a ambos un \omterm{def:g3:division:remainder}{cociente} \emph{y} a
\omterm{def:g3:division:remainder}{resto}.

\section{Escribir y comparar números enteros}

\begin{definition}[valor posicional]\label{def:g6:wholes:place}
Nuestra forma de escribir números utiliza diez dígitos ($0$ a $9$), y el
\emph{El valor de un dígito depende de su lugar.}\index{El valor de un dígito depende de su lugar.}: en
$5\,208$, el dígito $5$ cuenta miles, el $2$ cuenta cientos, el
$0$ cuenta decenas (no hay ninguna) y el $8$ cuenta unidades:
\[
5\,208 = 5 \times 1000 + 2 \times 100 + 0 \times 10 + 8 .
\]
\end{definition}

\begin{omfigure}
\begin{tikzpicture}[scale=0.9]
  \draw (0,0) grid[xstep=2.3, ystep=0.8] (9.2,1.6);
  \node[font=\small] at (1.15,1.2) {thousands};
  \node[font=\small] at (3.45,1.2) {hundreds};
  \node[font=\small] at (5.75,1.2) {tens};
  \node[font=\small] at (8.05,1.2) {units};
  \node[font=\large, omDef] at (1.15,0.4) {$5$};
  \node[font=\large, omDef] at (3.45,0.4) {$2$};
  \node[font=\large, omDef] at (5.75,0.4) {$0$};
  \node[font=\large, omDef] at (8.05,0.4) {$8$};
\end{tikzpicture}

{\small La tabla de valor posicional de $5\,208$. El $0$ importa: sin él,
el $5$ y el $2$ se deslizaría en las columnas equivocadas y el número
leería $528$.}
\end{omfigure}

\begin{method}[Comparar dos números enteros]\label{met:g6:wholes:compare}
\begin{enumerate}
  \item El número con más dígitos es el mayor
        ($1\,203 > 989$).
  \item Si tienen el mismo número de dígitos, compare dígito por dígito
        desde el \emph{izquierda}; la primera diferencia decide:
        $6\,742 > 6\,715$ porque en el lugar de las decenas, $4 > 1$.
\end{enumerate}
Los símbolos son $<$ ("menos que") y $>$ (``mayor que''); el
El extremo pequeño del símbolo apunta al número pequeño.
\end{method}

\begin{example}\label{ex:g6:wholes:order}
Ordene de menor a mayor: $908$; $89$; $1\,001$; $980$.
Primero por recuento de dígitos: $89$ (dos dígitos) viene primero, $1\,001$ (cuatro
dígitos) al final. Entre $908$ y $980$: mismo dígito de centenas $9$, entonces
decenas $0 < 8$, entonces $908 < 980$. Respuesta final:
\[
89 < 908 < 980 < 1\,001 .
\]
\end{example}

\begin{omfigure}
\begin{tikzpicture}[scale=1.05]
  \draw[-Stealth] (-0.4,0) -- (10.8,0);
  \foreach \x in {0,1,...,10} \draw (\x,-0.08) -- (\x,0.08);
  \foreach \x/\l in {0/0, 2/200, 4/400, 6/600, 8/800, 10/1000}
    \node[below, font=\small] at (\x,-0.1) {$\l $};
  \fill[omThm] (0.89,0) circle (2pt) node[above, font=\small] {$89$};
  \fill[omDef] (9.08,0) circle (2pt) node[above, font=\small] {$908$};
  \fill[omMeth] (9.8,0) circle (2pt) node[above right=-2pt, font=\small] {$980$};
\end{tikzpicture}

{\small Los números viven en una línea: más pequeño significa más a la izquierda. Cada tic
aquí vale la pena $100$.}
\end{omfigure}

\section{Suma, resta, multiplicación.}

\begin{example}[Cálculos de columnas]\label{ex:g6:wholes:columns}
Para sumar o restar, alinee las unidades debajo de las unidades, las decenas debajo de las
decenas, y trabaje de derecha a izquierda, llevando cuando sea necesario:
\[
\begin{array}{r}
4\,5\,7 \\
+\ 2\,8\,6 \\
\hline
7\,4\,3
\end{array}
\qquad\qquad
\begin{array}{r}
6\,0\,3 \\
-\ 1\,4\,7 \\
\hline
4\,5\,6
\end{array}
\]
Además: $7 + 6 = 13$, escribir $3$, llevar $1$; entonces
$5 + 8 + 1 = 14$, escribir $4$, llevar $1$; entonces $4 + 2 + 1 = 7$.
Comprueba la resta sumando de nuevo: $456 + 147 = 603$.
\end{example}

\begin{remark}[Vocabulario]
El resultado de una suma es \emph{suma}\index{suma}; de una resta,
un \emph{diferencia}\index{diferencia}; de una multiplicación, una
\emph{producto}\index{producto}; de una división, una
\emph{consciente}\index{consciente}. Los números que se combinan son los
\emph{términos} (para $+$, $-$) o \emph{factores} (para $\times $).
\end{remark}

\begin{example}[Computación inteligente]\label{ex:g6:wholes:smart}
Reordenar términos o factores a menudo ahorra trabajo:
\[
17 + 58 + 3 = (17 + 3) + 58 = 20 + 58 = 78 ,
\]
\[
25 \times 17 \times 4 = (25 \times 4) \times 17 = 100 \times 17 = 1700 .
\]
\end{example}

\section{División}

\begin{theorem}[división euclidiana]\label{thm:g6:wholes:division}
Dados dos números enteros $ a $ (el \emph{dividendo}) y $ b \neq 0$ (el
\emph{\omterm{def:g5:division:multiple}{divisor}}), hay exactamente una manera de escribir
\[
a = b \times q + r
\qquad\text{con}\qquad
0 \leq r < b .
\]
el número $ q $ es el \emph{consciente}\index{consciente} y $ r $ el
\emph{resto}\index{resto} de la división de $ a $ por $ b $.
\end{theorem}
\admitted

\begin{example}\label{ex:g6:wholes:division}
Dividir $137$ por $4$, paso a paso.
\begin{enumerate}
  \item ¿Cuántas veces $4$ entrar en $13$? tres veces
        ($4 \times 3 = 12$), \omterm{def:g3:division:remainder}{resto} $13 - 12 = 1$.
  \item Bajar el $7$: cuantas veces $4$ entrar en $17$?
        cuatro veces ($4 \times 4 = 16$), \omterm{def:g3:division:remainder}{resto} $1$.
\end{enumerate}
Entonces $ q = 34$ y $ r = 1$:
\[
137 = 4 \times 34 + 1, \qquad\text{and indeed } 0 \leq 1 < 4 .
\]
Controlar: $4 \times 34 = 136$, y $136 + 1 = 137$.
\end{example}

\begin{omfigure}
\begin{tikzpicture}[scale=0.95]
  \draw[-Stealth] (-0.3,0) -- (11.2,0);
  \foreach \i in {0,1,...,34} \draw ({\i*0.32},-0.06) -- ({\i*0.32},0.06);
  \foreach \i in {0,10,20,30}
    \node[below, font=\footnotesize] at ({\i*0.32},-0.08) {$4\times\i $};
  \draw[omDef, very thick] (0,0.14) -- (10.88,0.14);
  \draw[omThm, very thick] (10.88,0.14) -- (10.96,0.14);
  \fill (10.96,0) circle (1.8pt) node[above, font=\small] {$137$};
  \node[omDef, font=\small] at (5.4,0.55) {$34$ full groups of $4$};
\end{tikzpicture}

{\small Divisor $137$ por $4$: puedes encajar $34$ grupos completos de $4$
(alcanzando $136$), y $1$ sobra --- el \omterm{def:g3:division:remainder}{resto} es siempre
más pequeño que el \omterm{def:g5:division:multiple}{divisor}.}
\end{omfigure}

\begin{definition}[Divisible]\label{def:g6:wholes:divisible}
cuando el \omterm{def:g3:division:remainder}{resto} es $0$, decimos que $ a $ es
\emph{divisible}\index{divisible} por $ b $: por ejemplo $135$ es divisible por $5$ desde $135 = 5 \times 27 + 0$.
\end{definition}

\begin{method}[Problemas verbales con división]\label{met:g6:wholes:problems}
Después de computar $ a = bq + r $, siempre regrese a la pregunta:
\begin{enumerate}
  \item ``¿Cuántas casillas/equipos/páginas completas?'' --- la respuesta es la
        \omterm{def:g3:division:remainder}{cociente} $ q $;
  \item "¿Cuántos sobran?" --- la respuesta es la \omterm{def:g3:division:remainder}{resto}
        $ r $;
  \item "¿Cuántas cajas se necesitan para todos?" --- si $ r > 0$, uno
        \emph{más} que el \omterm{def:g3:division:remainder}{cociente}: $ q + 1$.
\end{enumerate}
\end{method}

\begin{example}\label{ex:g6:wholes:bus}
$137$ Los estudiantes realizan un viaje en autobuses de $40$ asientos. División:
$137 = 40 \times 3 + 17$. Tres autobuses transportan $120$ estudiantes, y $17$
los estudiantes se quedan --- ¡ellos también necesitan un autobús! Entonces $4$ Se necesitan autobuses, incluso
aunque el \omterm{def:g3:division:remainder}{cociente} es $3$.
\end{example}

\section{Ejercicios}

\begin{exercise}[$\star $]\label{exo:g6:wholes:1}
Escriba en dígitos: ``tres mil cuarenta y siete''; ``veinte mil
quinientos''; ``ciento cuatro''. Entonces escribe $60\,013$ en palabras.
\end{exercise}

\begin{exercise}[$\star $]\label{exo:g6:wholes:2}
en el número $47\,209$: ¿cuál es el dígito de las centenas? El dígito de
las unidades? ¿Qué significa el dígito? $4$ ¿contar? Escribe el número como \omterm{def:g1:addition:def}{suma} de
\omterm{def:g5:division:multiple}{múltiplos} de $1$, $10$, $100$, \dots\ como en
\cref{def:g6:wholes:place}.
\end{exercise}

\begin{exercise}[$\star $]\label{exo:g6:wholes:3}
Copiar y completar con $<$ o $>$:
\[
507 \;?\; 570, \qquad
1\,099 \;?\; 989, \qquad
23\,456 \;?\; 23\,465, \qquad
9\,999 \;?\; 10\,000 .
\]
\end{exercise}

\begin{exercise}[$\star $]\label{exo:g6:wholes:4}
Ordene de menor a mayor: $2\,043$; $978$; $2\,304$; $2\,034$;
$1\,001$.
\end{exercise}

\begin{exercise}[$\star $]\label{exo:g6:wholes:5}
Calcular en columnas: $348 + 576$; \quad $805 - 268$; \quad
$307 \times 26$. Comprueba la resta con una suma.
\end{exercise}

\begin{exercise}[$\star $]\label{exo:g6:wholes:6}
Calcula inteligentemente, agrupando términos o factores:
\[
38 + 45 + 12 + 5, \qquad
2 \times 83 \times 50, \qquad
4 \times 78 \times 25 .
\]
\end{exercise}

\begin{exercise}[$\star $]\label{exo:g6:wholes:7}
Para cada división, dé la \omterm{def:g3:division:remainder}{cociente} y el \omterm{def:g3:division:remainder}{resto}y escribe el
igualdad $ a = b \times q + r $:
\[
95 \div 7, \qquad
230 \div 6, \qquad
504 \div 8 .
\]
\end{exercise}

\begin{exercise}[$\star $]\label{exo:g6:wholes:8}
Sin dividir, explica por qué la igualdad. $173 = 8 \times 20 + 13$ es
\emph{no} la división euclidiana de $173$ por $8$ y encontrar el correcto
\omterm{def:g3:division:remainder}{cociente} y \omterm{def:g3:division:remainder}{resto}.
\end{exercise}

\begin{exercise}[$\star\star $]\label{exo:g6:wholes:9}
Los huevos se envasan en cajas de $12$.
\begin{enumerate}
  \item una granja tiene $2\,000$ huevos. ¿Cuántas cajas llenas puede llenar y
        ¿cuantos huevos sobran?
  \item Una panadería necesita $150$ huevos. ¿Cuántas cajas debe comprar?
\end{enumerate}
\end{exercise}

\begin{exercise}[$\star\star $]\label{exo:g6:wholes:10}
Hoy es martes. ¿En qué día de la semana será? $100$ ¿días?
(Las semanas tienen $7$ días: divide y piensa en el \omterm{def:g3:division:remainder}{resto}.)
\end{exercise}

\begin{exercise}[$\star\star $]\label{exo:g6:wholes:11}
en una división por $9$, el \omterm{def:g3:division:remainder}{cociente} es $56$ y el \omterm{def:g3:division:remainder}{resto} es el
más grande que pueda ser. ¿Qué es el dividendo?
\end{exercise}

\begin{exercise}[$\star\star\star $]\label{exo:g6:wholes:12}
Soy un número de tres dígitos. Mi dígito de centenas es el doble de mi dígito de unidades,
mi dígito de las decenas es $0$, y el \omterm{def:g1:addition:def}{suma} de mis dígitos es $9$. ¿Quién soy yo?
(Razone paso a paso y luego verifique).
\end{exercise}

\section{Problema: el pequeño Gauss y la suma de las primeras cien
números}

\begin{problem}[{Problema de fin de semana --- sumando $1 + 2 + 3 + \dots + 100$ en diez segundos, y los números triangulares}]
\label{pb:g6:wholes:1}
Una historia famosa: para mantener ocupada a su clase, un maestro de escuela pidió a su
los alumnos sumen todos los números de $1$ a $100$. un niño escribió
$5\,050$ en su pizarra casi de inmediato --- Carl Friedrich Gauss, nueve
años, que se convertiría en uno de los más grandes matemáticos
de todos los tiempos. Su secreto no era más que \emph{computación inteligente}
(\cref{ex:g6:wholes:smart}): agrupar los términos para que cada grupo
es fácil. Este problema redescubre su truco, lo hace funcionar en cada
situación, y muestra con qué frecuencia estos \omterm{def:g1:addition:def}{sumas} aparecen en la vida diaria.

\medskip
\noindent\textbf{Parte I --- Gauss's pairs.}
\begin{enumerate}
  \item Calcular $1 + 2 + 3 + 4 + 5$ agrupando términos inteligentemente.
  \item Para $1 + 2 + 3 + \dots + 10$: empareja el primer término con el
        último ($1 + 10$), el segundo con el penúltimo
        ($2 + 9$), etcétera. ¿Cuánto vale cada par? cuantos
        ¿Hay pares? Termina el cálculo.
  \item Mismo método para $1 + 2 + \dots + 20$: ¿Qué es cada par?
        vale, ¿cuántos pares, qué total?
  \item Ahora el maestro de escuela \omterm{def:g1:addition:def}{suma}, $1 + 2 + \dots + 100$: explicar
        por que las parejas $1 + 100$, $2 + 99$, $3 + 98, \dots $ son todos
        valen lo mismo, y por qué hay exactamente $50$ de ellos,
        siendo el ultimo $50 + 51$.
  \item Escribe la respuesta de Gauss como una única multiplicación y
        calcularlo.
\end{enumerate}

\medskip
\noindent\textbf{Parte II --- A trick that never fails.}
\begin{enumerate}[resume]
  \item Pruebe el emparejamiento $1 + 2 + \dots + 9$: esta vez uno
        El número se queda sin pareja. ¿Cuál? Emparejar alrededor
        y calcular el \omterm{def:g1:addition:def}{suma}.
  \item Aquí hay una versión del truco que funciona ya sea que el
        el recuento de términos es \omterm{def:g3:numbers:evenodd}{incluso} o \omterm{def:g3:numbers:evenodd}{extraño}: escribe el \omterm{def:g1:addition:def}{suma} dos veces, una vez
        hacia adelante y una vez hacia atrás, una línea debajo de la otra,
        \[
        \begin{array}{ccccccc}
        1 & + & 2 & + & \dots & + & 9 \\
        9 & + & 8 & + & \dots & + & 1
        \end{array}
        \]
        y agregar \emph{columna por columna}. ¿Qué es cada columna?
        ¿Vale la pena? ¿Cuántas columnas hay? Explique por qué el \omterm{def:g1:addition:def}{suma} de
        las dos lineas son $9 \times 10 = 90$, y concluir que
        $1 + 2 + \dots + 9 = 45$ --- la misma respuesta que en
        pregunta 6.
  \item Usa el truco de las dos líneas para calcular $1 + 2 + \dots + 50$,
        entonces $1 + 2 + \dots + 200$.
  \item Calcular el \omterm{def:g1:addition:def}{suma} de los primeros cien \emph{\omterm{def:g3:numbers:evenodd}{incluso}} números,
        $2 + 4 + 6 + \dots + 200$. (Compare cada término con un término
        de Gauss \omterm{def:g1:addition:def}{suma}: no se necesita un nuevo emparejamiento.)
  \item Calcular $101 + 102 + \dots + 200$, usando dos \omterm{def:g1:addition:def}{sumas} tu
        ya lo sé y una resta.
\end{enumerate}

\medskip
\noindent\textbf{Parte III --- \omterm{pb:g6:wholes:1}{Números triangulares} everywhere.}
los totales $1$, $3$, $6$, $10$, $15, \dots $ del \omterm{def:g1:addition:def}{sumas}
$1 + 2 + \dots + n $ son llamados \emph{triangular
números}\index{triangular
números}, porque cuentan objetos
apilados en un triángulo.
\begin{enumerate}[resume]
  \item Ocho amigos se encuentran y cada uno se da la mano exactamente una vez.
        con cada uno de los demás. Explique por qué el número de
        apretones de manos es $7 + 6 + 5 + 4 + 3 + 2 + 1$ y calcularlo.
  \item Un tendero construye un triángulo. \omterm{def:g5:solids:def}{pirámide} de latas: $12$ latas en
        la fila inferior, $11$ en el siguiente, y así sucesivamente hasta un solo
        lata encima. ¿Cuántas latas necesita? (Usa el truco de
        pregunta 7.)
  \item Dibuja el \omterm{def:g1:addition:def}{suma} $1 + 2 + 3 + 4$ como una escalera de plazas:
        un cuadrado en la primera columna, dos en la segunda, tres
        en el tercero, cuatro en el cuarto. Mostrar en tu dibujo
        eso \emph{dos} tales escaleras, una de ellas volteada
        abajo, encajan exactamente en un rectángulo de $4$ filas de
        $5$ cuadrados --- y explica por qué este es el truco de las dos líneas
        de la pregunta 7, dibujada en lugar de escrita.
  \item suena un reloj $1$ hora a la una, $2$ veces a las dos
        en punto, \dots, $12$ veces a las doce en punto. cuantos
        huelgas es que durante las doce horas? Y sobre un todo
        día?
  \item Una última para Gauss: $5\,050$ segundos --- ¿eso es más o menos?
        menos de una hora y un \omterm{def:g3:division:half}{medio}? Conviértelo en horas,
        minutos y segundos, usando la división euclidiana
        (\cref{thm:g6:wholes:division}) dos veces.
\end{enumerate}
\end{problem}
